
% ============================================================================
\clearpage
\section*{v4.8.0-r09 Source-Only: Tests of P2 and P3, and Re-Classification of Earlier Results}
\addcontentsline{toc}{section}{v4.8.0-r09 Source-Only: Tests of P2 and P3, and Re-Classification of Earlier Results}
\label{sec:v480-r09-p2-p3-reclassification}
% ============================================================================

\paragraph{Purpose.}
Corollary~\ref{cor:falsifiable} lists P2 (Koide phase) as ``already matches data'' and P3
(\(\bsym\)-formula) as ``matches to \(0.018\%\)''.  This layer tests both against the measurement
precision and against the freedom in choosing the formula.  Since layers r03--r08 changed the status
of many earlier statements, it then collects the current status of all predictions and of the main
results in one register, so that no earlier ``proved'' or ``\checkmark'' is read without its
present qualification.  Earlier layers are not rewritten.

\noindent\textbf{Source-only status.}
This v4.8.0-r09 layer is source-only.  It does not build a PDF, does not assign a new DOI,
and does not authorize a public release.  The v4.8.0 Version DOI
\href{https://doi.org/10.5281/zenodo.23365244}{10.5281/zenodo.23365244} and the Concept DOI
\href{https://doi.org/10.5281/zenodo.19210728}{10.5281/zenodo.19210728} are unchanged by it.

\subsection*{v4.8.0-r09 P2: Koide phase}
\label{subsec:v480-r09-p2}

Theorem~\ref{thm:koide-triolic} predicts, in the Brannen form
\(\sqrt{m_k}=A\bigl(1+\sqrt2\cos(\theta_K+2\pi k/3)\bigr)\), the phase
\(\theta_K=\pi-\sqrt2\arcsin\bsym=2.319109\) rad \(=132.8751^\circ\).

\paragraph{Theoremlet v4.8.0-r09.1.}
From \(m_e=0.51099895\), \(m_\mu=105.6583755\) and \(m_\tau=1776.93\pm0.09\) MeV the observed phase is
\(\theta_K^{\rm obs}=132.7325^\circ\pm0.0004^\circ\) (with \(m_\tau=1776.86\pm0.12\):
\(132.7328^\circ\pm0.0005^\circ\)).  The prediction differs by \(+0.142^\circ\), i.e.\ by roughly
\(300\)--\(400\) standard deviations.  P2 is excluded by the data.

\paragraph{Proof sketch.}
Solve the Brannen form for \(\theta\) from the three masses; propagate the \(\tau\)-mass error, which
dominates.  Script \texttt{k12\_p2\_p3.py} in the working folder.  \(\square\)

\paragraph{Remarks.}
(i) The relative deviation of \(0.11\%\) quoted in Theorem~\ref{thm:koide-triolic} is small compared with
the angle, but large compared with the measurement error; a prediction is tested against the error.
(ii) The observed phase is \(\theta_K^{\rm obs}=2\pi/3+0.22223\), within \(10^{-5}\) of
\(2\pi/3+2/9\); this empirical regularity is known in the literature and is not produced by the formula
of Theorem~\ref{thm:koide-triolic}.
(iii) Step~1 of the proof of Theorem~\ref{thm:koide-triolic} derives \(K=2/3\) from Cauchy--Schwarz.
Cauchy--Schwarz gives \((\sum\sqrt{m_i})^2\le3\sum m_i\), i.e.\ \(K\ge1/3\), with equality only for equal
masses; it does not yield \(K=2/3\).  The Koide value \(K=0.666665\) is an empirical fact; its derivation
is open.

\subsection*{v4.8.0-r09 P3: the \texorpdfstring{$\bsym$}{beta}-formula}
\label{subsec:v480-r09-p3}

Remark~\ref{conj:cl-correction} states \(\bsym=\tfrac1{\sqrt3}\exp\bigl(-\gamma_1/(4\pi\gL^5)\bigr)\).

\paragraph{Theoremlet v4.8.0-r09.2.}
With \(\gL=\gBI/\gstring=1.86444\) the formula gives \(0.549233\) against
\(\bsym=\sqrt{(\gL-1)/(\gL+1)}=0.549348\), a deviation of \(-0.021\%\); it holds exactly for
\(\gL=1.86366\).  It is a relation between the two inputs \(\gBI,\gstring\) and \(\gamma_1\), not a
prediction of an independently measured quantity.  Among the 56 formulas
\(\tfrac1{\sqrt3}\exp(-\gamma_1/(c\,\gL^n))\), \(n=1,\dots,8\),
\(c\in\{1,2,\pi,2\pi,4\pi,8\pi,16\pi\}\), 27 lie within \(\pm5\%\); the chance that one of these lands
within \(\pm0.02\%\) is about \(11\%\).

\paragraph{Proof sketch.}
Both sides are functions of \(\gL\) alone; solve by bisection.  Enumerate the family and estimate the
look-elsewhere probability with a uniform density over \(\pm5\%\).  \(\square\)

\paragraph{Assessment.}
The exponent was changed from \(\gL^2\) to \(\gL^5\) after the fact (Remark~\ref{conj:cl-correction}),
and Remark~\ref{conj:cl-correction} itself marks it as requiring structural derivation.  P3 is
recorded as a numerical coincidence of moderate significance, not as a confirmed prediction.  A
derivation of the exponent before comparison with data would change this.

\subsection*{v4.8.0-r09 register of predictions}
\label{subsec:v480-r09-predictions}

\begin{center}
\small
\begin{tabular}{@{}>{\raggedright\arraybackslash}p{0.27\textwidth}>{\raggedright\arraybackslash}p{0.31\textwidth}>{\raggedright\arraybackslash}p{0.34\textwidth}@{}}
\toprule
prediction (Cor.~\ref{cor:falsifiable}) & earlier status & status after r09\\
\midrule
P1 GUE photon-frame unfolding & testable now & not confirmed (r08); minimum at \(\mu=1\)\\
P2 Koide phase & matches to \(0.111\%\) & excluded by data, \(\sim300\)--\(400\sigma\) (r09.1)\\
P3 \(\bsym\)-formula & matches to \(0.018\%\) & coincidence, exponent chosen post hoc; \(\sim11\%\) look-elsewhere (r09.2)\\
P4 \(\Lambda\) as frame artefact & reinterpretation & contradicted as stated (r03 F4); \(\Lambda\) type open (r03 F3); \(\Theta\)-sum route false (r05)\\
P5 \(\alpha\)-drift ratio & not yet testable & exponent corrected (r03 F2); untestable at present\\
P6 Hawking, tachyon frame & theoretical & assumption \(k(M)=0\) made explicit (r03 F6); untested\\
\bottomrule
\end{tabular}
\end{center}

\subsection*{v4.8.0-r09 register of main results}
\label{subsec:v480-r09-results}

\begin{center}
\small
\begin{tabular}{@{}>{\raggedright\arraybackslash}p{0.34\textwidth}>{\raggedright\arraybackslash}p{0.58\textwidth}@{}}
\toprule
statement & current status (layer)\\
\midrule
Thm.~\ref{thm:unity}, equivalence of (i)--(v) & conjecture; Stone--von Neumann justification withdrawn (r04); Reading~B gives a precise form via duality (r07)\\
Thm.~\ref{thm:unitary-equiv} & open (r04)\\
Def.~\ref{def:biquaternion}, Minkowski reading & replaced: signature \((1,6)\); natural norm \(N(q)\), Reading~B (r06)\\
Thm.~\ref{thm:lorentz} & shows \(J\)-invariance; Lorentz action from Reading~B (r06)\\
Thm.~\ref{thm:hopf} & Hopf fibrations classical; physical identification conjecture; the chain \(S^{15}\to S^8\to S^4\to S^2\) is not a composition of Hopf maps (recorded here)\\
Thm.~\ref{thm:fractal} & null axis redefined as \(s=0\), directions \(=S_0\) (r06); transport open\\
triole structure & triole lemma proved (r06.2); octonionic triole configuration repaired (r04.4)\\
Heisenberg compensator, cascade & standalone, proved (r04.2, r04.3)\\
Thm.~\ref{thm:einstein} & structural conjecture; action undefined (r05)\\
Thm.~\ref{thm:gue} & empirical fact; derivation withdrawn (r05)\\
Thm.~\ref{thm:gue-selfdet} & not confirmed (r08)\\
Thm.~\ref{thm:koide-triolic} & \(K=2/3\) empirical, derivation open; phase formula excluded (r09)\\
Thm.~\ref{thm:lambda-zero} & (i) holds by definition of the photon frame; (iii) value off by \(10^{121}\) (r03 F8, F9)\\
Cor.~\ref{cor:matter-error} & corrected to \(\gL^5\approx22.5\) (r03 F1)\\
\(\gL=\gBI/\gstring\) & ratio of two published Barbero--Immirzi values (r03 F11); input, not derived (OP1)\\
OP5 Step~(B) & no valid form (r05)\\
RH & not proved; forbidden as global claim (\(W_\gamma^{\rm glob}\)); Riemann source no proof source (r04)\\
\bottomrule
\end{tabular}
\end{center}

\paragraph{Theoremlet v4.8.0-r09.3.}
After r09 no prediction of Corollary~\ref{cor:falsifiable} stands as confirmed.  The results that carry
are the structural ones: the standalone compensator and cascade, the repaired triole hierarchy, the
triole lemma, the signature analysis and Reading~B with its orbit, null-cone and duality structure.
None of the five guards is affected.

\paragraph{Proof sketch.}
Collect the decisions of r03--r09 as listed in the two registers.  \(\square\)

\subsection*{v4.8.0-r09 ledger and guard update}
\label{subsec:v480-r09-ledger-guard}

\[
R_{\rm pred}^{v480\text{-}r09}=(0,0,0,1)
\]
for (P2 decided; P3 classified; register compiled; derivations of \(K=2/3\) and of the P3 exponent open).

The layer spends no guard:
\[
R_{\rm glob}^{(v480\text{-}r09,+)}=(0,1,1,1,1),
\qquad
\Gamma_{\rm guard}^{v480\text{-}r09}=(1,1,1,1,1),
\qquad
\beta_4=1/256\ \text{unspent}.
\]
It does not imply exact finite \(C5\) zero closure, exact all-readout terminality,
\(C6/T2\), a global photon-clock/RH theorem, or global all-readout closure.

